
% ============================================================
% !TeX root = ../main.tex
% Resumo (Portuguese abstract — MANDATORY at ISCTE)
% ============================================================
\chapter*{Resumo}

\selectlanguage{portuguese}
\noindent
A análise da geometria e da espessura da parede do ventrículo esquerdo
é relevante para a caracterização da estrutura e função cardíacas.
Contudo, a ressonância magnética cardíaca em eixo curto fornece
uma representação bidimensional através de cortes espacialmente
separados, tornando necessária a reconstrução tridimensional da
geometria entre os cortes disponíveis. Embora a interpolação permita
estabelecer continuidade entre contornos, não garante a preservação
da anatomia ventricular. Além disso, a avaliação clínica da espessura
miocárdica baseia-se frequentemente em medições localizadas e na
análise visual. Neste contexto, o modelo de 17 segmentos da American
Heart Association (AHA) constitui uma referência para a avaliação
regional do ventrículo esquerdo.

Nesta dissertação, propõe-se uma abordagem baseada em aprendizagem
automática para a reconstrução tridimensional do ventrículo esquerdo
a partir de contornos de ressonância magnética cardíaca. Utiliza-se
uma representação neuronal implícita para descrever continuamente
as superfícies endocárdica e epicárdica nas fases de telediástole
e telessístole. O modelo é inicialmente treinado com formas
ventriculares sintéticas e posteriormente ajustado com segmentações
clínicas. A espessura da parede miocárdica é estimada a partir da
geometria reconstruída, permitindo a análise local e regional segundo
os 17 segmentos da AHA.

O desempenho da abordagem é avaliado por comparação com métodos
baseados em funções de base radial e com um modelo estatístico
de forma. Os resultados mostram que as superfícies reconstruídas
acompanham os contornos observados, com maior concordância entre
métodos na cavidade ventricular do que na parede miocárdica.
Verifica-se também que as estimativas da espessura dependem da
geometria reconstruída e do método de cálculo, pelo que uma correlação
elevada não implica necessariamente concordância entre medições.
A análise segmentar permite ainda caracterizar as variações regionais
da espessura entre a telediástole e a telessístole.

Este trabalho explora, assim, o potencial das representações neuronais
implícitas na reconstrução tridimensional do ventrículo esquerdo e
na quantificação da espessura miocárdica. Contudo, na ausência de
uma referência anatómica tridimensional independente, não é possível
determinar conclusivamente qual dos métodos produz a reconstrução
mais próxima da anatomia real. Os resultados evidenciam, ainda assim,
a influência das opções de modelação geométrica nas estimativas
da espessura e a importância de considerar o método de reconstrução
na interpretação das medições.

\vspace{1em}
\noindent\textbf{Palavras-chave:} Ressonância magnética cardíaca;
ventrículo esquerdo; reconstrução tridimensional; representações
neuronais implícitas; funções de distância com sinal; espessura
da parede miocárdica.

\selectlanguage{english}
\clearpage
